\documentclass[11pt]{article}
\usepackage{geometry}
\geometry{margin=1in}
\usepackage{amsmath}
\usepackage{enumitem}
\usepackage{booktabs}

\title{STAT115 2024 Practice Exam Solutions}
\author{Compiled via Codex}
\date{}

\begin{document}
\maketitle

\begin{enumerate}[label=\textbf{Q\arabic*:}, leftmargin=*, itemsep=0.9em]
\item \textbf{Answer: B.} Compute the sample mean: $(36 + 49.1 + 41.5 + 40 + 46.4)/5 = 42.6\,\text{kg}$.
\item \textbf{Answer: E.} Use the sample standard deviation with denominator $n-1$: $s = \sqrt{\frac{\sum (x_i - \bar{x})^2}{4}} \approx 5.20\,\text{kg}$.
\item \textbf{Answer: C.} The measurements are quantitative on a continuum, so the data type is continuous.
\item \textbf{Answer: A.} The $65$ minutes is the observed sample mean from $n=764$ students, so it is a statistic and a realised value of a random variable.
\item \textbf{Answer: C.} The target population is all New Zealand high school students and the parameter is their mean weekly exercise time.
\item \textbf{Answer: B.} Treating 1989--2000 as 12 years, the incidence rate is $\frac{175}{1904\times 12}\times 100{,}000 \approx 7.66\times 10^2$ fractures per $100{,}000$ person-years, describing the study cohort.
\item \textbf{Answer: B.} The relative risk is $\big(175/(1904\times 12)\big)\big/\big(54/(1373\times 12)\big) \approx 2.34$.
\item \textbf{Answer: D.} Under-diagnosis of men makes the observed male rate too small, so the ratio $\text{rate}_\text{women}/\text{rate}_\text{men}$ is biased upward.
\item \textbf{Answer: A.} The quoted 4.5 and 1.9 values summarise the observed rate ratios in the study for the 60--69 and $80+$ age bands.
\item \textbf{Answer: D.} The mean of the total is $\mu_W = \mu_X + \mu_Y = 1750 + 750 = 2500$.
\item \textbf{Answer: C.} With independent counts, variances add: $\sigma_W = \sqrt{100^2 + 40^2} = \sqrt{11{,}600} \approx 1.08\times 10^2$.
\item \textbf{Answer: C.} $\Pr(B) = 82/623 \approx 0.1316$.
\item \textbf{Answer: A.} $\Pr(A^c \cap B) = 50/623 \approx 0.0803$.
\item \textbf{Answer: E.} $\Pr(A^c \mid B) = 50/82 \approx 0.610$.
\item \textbf{Answer: B.} Sensitivity is $\Pr(T\mid B) = 0.97$.
\item \textbf{Answer: A.} Specificity is $\Pr(T^c\mid B^c) = 1 - 0.35 = 0.65$.
\item \textbf{Answer: B.} $\Pr(B\cap T) = \Pr(B)\Pr(T\mid B) = 0.20\times 0.97 = 0.194$.
\item \textbf{Answer: C.} $\Pr(B^c \cap T) = 0.80\times 0.35 = 0.28$.
\item \textbf{Answer: A.} $\Pr(T) = 0.194 + 0.28 = 0.474 \approx 0.47$.
\item \textbf{Answer: C.} Positive predictive value $= 0.194/0.474 \approx 0.41$.
\item \textbf{Answer: B.} A binomial model requires a fixed $n$, identical success probability for each trial, independence, and two possible outcomes per trial.
\item \textbf{Answer: C.} From the table, $\Pr(X=3) = 0.0081$.
\item \textbf{Answer: D.} $\Pr(X>3) = \Pr(4) + \Pr(5) = 0.00045 + 0.00001 = 0.00046$.
\item \textbf{Answer: D.} Because 150 exceeds the mean of a symmetric normal distribution, the upper-tail probability must be between $0$ and $0.5$.
\item \textbf{Answer: A.} The upper tail is $1 - \operatorname{pnorm}(137.5,\ 127.5,\ 19.6)$.
\item \textbf{Answer: B.} The middle probability is $\operatorname{pnorm}(134.5, 127.5, 19.6) - \operatorname{pnorm}(119.5, 127.5, 19.6)$.
\item \textbf{Answer: D.} $z = (140.5 - 127.5)/19.6 \approx 0.66$.
\item \textbf{Answer: B.} Using $t_{0.975,194}$ with $s/\sqrt{n}$ constructs a 95\% confidence interval for the population mean BMI.
\item \textbf{Answer: D.} Using $z_{0.975}\,s$ without $\sqrt{n}$ gives an estimated 95\% reference range for individual BMIs.
\item \textbf{Answer: C.} The mean paired difference (after $-$ before) is $52.9 - 56.5 = -3.6$ bpm.
\item \textbf{Answer: E.} The estimated standard error is $s_d/\sqrt{n} = 5.9/\sqrt{13} \approx 1.64$ bpm.
\item \textbf{Answer: A.} For a 95\% interval with $n=13$, use $t_{0.975,12}$.
\item \textbf{Answer: D.} The 95\% CI for the population mean difference excludes zero and is entirely negative, so the data suggest a true decrease in resting heart rate among female runners following the protocol.
\item \textbf{Answer: A.} The sample proportion is $\hat{p} = 49/99 \approx 0.495$.
\item \textbf{Answer: B.} $\text{SE} = \sqrt{\hat{p}(1-\hat{p})/n} = \sqrt{0.495\times 0.505/99} \approx 0.050$.
\item \textbf{Answer: B.} A large-sample CI for a proportion uses $\hat{p} \pm z_{0.975}\,\text{SE}$.
\item \textbf{Answer: C.} The interval gives a 95\% confidence range for the population proportion of New Zealand adults whose main news source is mainstream media.
\item \textbf{Answer: E.} A valid normal approximation requires a large sample and success probability not too close to 0 or 1; $\hat{p}\approx 0.5$ satisfies this.
\item \textbf{Answer: B.} If mainstream-media users respond more, the sample proportion is inflated—a form of selection bias upward.
\item \textbf{Answer: C.} Two independent cross-sectional surveys describe prevalence; no exposure is manipulated, so this is descriptive.
\item \textbf{Answer: A.} $\hat{p}_2 - \hat{p}_1 = 110/990 - 180/970 \approx -0.074$.
\item \textbf{Answer: D.} $\text{SE} = \sqrt{\hat{p}_1(1-\hat{p}_1)/970 + \hat{p}_2(1-\hat{p}_2)/990}$.
\item \textbf{Answer: D.} A 99\% two-sided CI uses $z_{0.995}$: $(\hat{p}_2 - \hat{p}_1) \pm z_{0.995}\,\text{SE}$.
\item \textbf{Answer: A.} The entire 99\% CI is negative, supporting a reduction in the proportion of daily smokers between 2010 and 2020.
\item \textbf{Answer: C.} Participants were randomised to intervention or control, so this is an experimental analytic randomised controlled trial.
\item \textbf{Answer: C.} To compare population means: $H_0\!:\mu_I = \mu_C$ versus $H_A\!:\mu_I \ne \mu_C$.
\item \textbf{Answer: B.} $\text{SE} = \sqrt{24/102 + 12.5/91} \approx 0.61$ units of reduction.
\item \textbf{Answer: A.} Test statistic $t = (9.6 - 8.3)/0.61 \approx 2.13$.
\item \textbf{Answer: D.} For a two-sided $t$-test use $2\,(1 - \operatorname{pt}(|t|,\nu))$.
\item \textbf{Answer: C.} With $p=0.0345 < 0.05$ we reject $H_0$ and conclude there is evidence of a difference in mean mercury reduction between groups.
\item \textbf{Answer: D.} The $p$-value is the probability, under $H_0$, of observing a difference at least as extreme as 1.3 units (in absolute value).
\item \textbf{Answer: C.} Increasing both the sample sizes and the significance level (e.g. $\alpha=0.1$) increases statistical power.
\item \textbf{Answer: B.} Farmer odds $= 53/447 \approx 0.119$.
\item \textbf{Answer: D.} Non-farmer odds $= 27/473 \approx 0.057$.
\item \textbf{Answer: A.} Odds ratio $= (53/447)/(27/473) = (53\times 473)/(447\times 27) \approx 2.08$.
\item \textbf{Answer: E.} $\text{SE}_{\log \text{OR}} = \sqrt{1/53 + 1/27 + 1/447 + 1/473} \approx 0.2455$.
\item \textbf{Answer: B.} Exponentiating $(0.2499, 1.2121)$ yields $(e^{0.2499}, e^{1.2121}) \approx (1.28, 3.36)$.
\item \textbf{Answer: A.} For an odds ratio, bounds entirely above/below 1 indicate increased/decreased odds; intervals spanning 1 imply insufficient evidence.
\item \textbf{Answer: A.} Missed diagnoses among farmers reduce observed cases, biasing the OR downward; this is information bias.
\item \textbf{Answer: C.} If farmers are more often male and males have lower MDD risk, sex confounding would pull the crude odds ratio toward 1, masking part of the association.
\item \textbf{Answer: D.} The adjusted OR of 1.37 has a 95\% CI that includes 1, so after adjustment the data do not show a statistically significant association.
\item \textbf{Answer: E.} $H_0$: diagnosis and paternal age are independent; $H_A$: there is an association.
\item \textbf{Answer: A.} Expected count $= (20\times 528)/1612 \approx 6.55$ children.
\item \textbf{Answer: B.} Degrees of freedom $(r-1)(c-1) = (2-1)(3-1) = 2$.
\item \textbf{Answer: C.} Right-tail $p$-value is $1 - \operatorname{pchisq}(13.30,\nu)$.
\item \textbf{Answer: A.} With $p=0.0014 < 0.05$ we reject $H_0$ and infer an association between paternal age and autism diagnosis.
\item \textbf{Answer: C.} Retinol level is modelled as the response variable and age as the explanatory variable.
\item \textbf{Answer: A.} The fitted slope is about $3.07$ ng/mL per year, as shown in the regression output.
\item \textbf{Answer: C.} Test $H_0\!:\beta_1 = 0$ versus $H_A\!:\beta_1 \ne 0$ for the age slope.
\item \textbf{Answer: C.} For simple regression with $n=100$, the $t$-statistic uses $\nu = n-2 = 98$ degrees of freedom.
\item \textbf{Answer: E.} The $p$-value $0.0172 < 0.05$ provides evidence of an association between age and retinol levels.
\item \textbf{Answer: E.} The intercept estimates the mean retinol level when age $=0$, i.e., for a newborn.
\item \textbf{Answer: B.} $\hat{\beta}_1$ estimates the mean change in retinol level for each additional year of age.
\item \textbf{Answer: C.} $R^2 = 0.05654$ indicates that about $5.654\%$ of the variance in retinol levels is explained by age.
\item \textbf{Answer: D.} Predicted value: $441.86 + 3.07\times 50 \approx 595.36$ ng/mL.
\item \textbf{Answer: B.} A 95\% prediction interval gives the range where a new 40-year-old's retinol level would fall with probability 0.95.
\item \textbf{Answer: D.} Prediction errors grow as age moves away from the sample mean (50.61 years), so forecasts at 40.61 years are less precise.
\item \textbf{Answer: B.} Predicted for age $=42$: $441.86 + 3.07\times 42 = 570.8$; residual $= 580.7 - 570.8 \approx 9.9$ ng/mL.
\item \textbf{Answer: A.} To model a binary outcome (cancer) with a continuous predictor, use logistic regression with cancer as the response.
\item \textbf{Answer: E.} The model predicts medical costs (charges), so charges is the response variable.
\item \textbf{Answer: B.} $\hat{y} = -12052.46 + 257.73\times 50 -128.64\times 1 + 322.36\times 28 + 474.41\times 3 + 23823.39\times 0 \approx 1.12\times 10^4$ dollars.
\item \textbf{Answer: B.} $\hat{\beta}_{\text{bmi}}$ estimates the change in expected charges for each one-unit BMI increase while holding the other predictors fixed.
\item \textbf{Answer: A.} The CI for the age slope is $257.73 \pm t_{0.975,1332}\times 11.9$.
\item \textbf{Answer: C.} The five predictors jointly explain about $74.97\%$ of the variation in medical charges.
\item \textbf{Answer: C.} $\mu_i$ denotes the population mean reduction for pesticide $i$, and $e_{ij}$ is the residual for plot $j$ using pesticide $i$.
\item \textbf{Answer: A.} ANOVA assumes $e_{ij}$ are independent $N(0,\sigma^2)$ within each group (common variance and zero mean).
\item \textbf{Answer: C.} $H_0$: $\mu_A=\mu_B=\mu_C=\mu_D=\mu_E=\mu_F$ versus the alternative that not all population means are equal.
\item \textbf{Answer: B.} From the ANOVA table, total SS $=2668.8+1015.2=3684.0$ and residual SS $=1015.2$.
\item \textbf{Answer: D.} Large between-group variability relative to within-group noise favours rejecting $H_0$ in one-way ANOVA.
\item \textbf{Answer: C.} The F-test $p$-value is $\operatorname{pf}(f,5,66,\text{lower.tail}=\text{FALSE})$.
\item \textbf{Answer: C.} The p-value $<0.01$ leads to rejecting $H_0$, indicating different mean reductions among at least some pesticides.
\item \textbf{Answer: B.} Blocking on country would absorb part of the residual variation, improving sensitivity to detect pesticide effects.
\end{enumerate}
\end{document}
